% ==============================================================================
% Executive Summary - Compact Front-Matter Layout
% ==============================================================================

\section*{Executive Summary}
\addcontentsline{toc}{section}{Executive Summary}

Switzerland occupies a disproportionately large position in global finance, consistently ranking among the world's top ten direct investment economies. This monograph investigates the structural trends, determinants, and macroeconomic spillovers of Swiss direct investment over the decade \textbf{2015--2024}, utilizing official \textbf{Swiss National Bank (SNB)} datasets.

\begin{findingbox}[Core Empirical Takeaways (2015--2024)]
\begin{itemize}[leftmargin=1.2em, itemsep=1.5pt]
    \item \textbf{Structural Capital Export Asymmetry:} Outward FDI stock ($\CHF\,1,340.1\bn$, $167.5\%$ GDP) exceeded inward FDI stock ($\CHF\,925.7\bn$, $115.7\%$ GDP) by $\CHF\,414.4\bn$ in 2024, confirming Switzerland's enduring net capital exporter role.
    \item \textbf{Post-2017 Inward Contraction:} Inward FDI contracted by $31.3\%$ from its 2017 peak ($\CHF\,1,347.8\bn$) to $\CHF\,925.7\bn$ in 2024, driven by OECD BEPS compliance, corporate tax reform (STAF/TRAF 2020), and Franc valuation deflation.
    \item \textbf{Phantom Conduit Distortion Unmasked:} While the Netherlands ($26.4\%$) and Luxembourg ($15.2\%$) lead immediate statistics, Ultimate Beneficial Owner (UBO) tracing unmasks the \textbf{United States as the dominant investor}, controlling $40.0\%$ ($\CHF\,370.7\bn$) of inward FDI.
    \item \textbf{Sectoral Asymmetry \& Global Labor Multiplier:} Inward FDI is skewed toward services ($77.6\%$), while outward FDI anchors high-value manufacturing ($37.9\%$, led by life sciences). Swiss MNEs operate $21,265$ foreign affiliates employing $2.48\mn$ personnel abroad ($4.43\times$ domestic parent employment).
\end{itemize}
\end{findingbox}
